% Insert after the current Waterbirds appendix and before
% J_forward_only_attribution_results.tex.

\section{Covertype Experimental Details and Complete Results}
\label{app:covertype-transfer}

This appendix reports the complete Covertype audit behind Section~\ref{sec:controlled}. The experiment uses a natural tabular base and controlled context--factor channels. Its purpose is not to assume that every treatment succeeds. It measures which response mechanism each trained model actually realizes.

\subsection{Setup and Operational Outcomes}
\label{app:covertype-setup}

We balance classes 1 and 2 from Covertype and retain 240,000 examples. The 54 natural features are augmented by one binary context and one binary candidate factor. The direction module uses context $G$ and factor $Z$; the fragility module uses context $H$ and factor $U$. Natural features and data splits are shared across all mechanism variants.

\begin{table}[t]
\centering
\caption{Covertype benchmark design. Each model family uses three seeds. Direction experiments additionally use shortcut strengths $p\in\{0.75,0.95\}$.}
\label{tab:covertype-setup}
\small
\begin{tabular}{llllr}
\toprule
Module & Regimes & Held-out operational behavior & Families & Models \\
\midrule
Direction & Direct, Gate, Invert & preserve / collapse / invert & 5 & 90 \\
Fragility & Robust, Mild, Fragile & endpoint-null prediction change & 5 & 45 \\
\midrule
Total & & & & 135 \\
\bottomrule
\end{tabular}
\end{table}

For the direction module, the endpoint is $G=+1$. We query the factor effect under $G=+1$ and $G=-1$, then define preservation, collapse, and inversion from the held-out signed responses. For the fragility module, the endpoint is $H=+1$. The primary behavioral target used in Section~\ref{sec:controlled} is the prediction-change rate under $H=-1$ among examples satisfying the endpoint-null gate. This target was stored by the formal experiment and matches the definition of $F$.

\subsection{Complete Behavior Alignment}
\label{app:covertype-complete-ranks}

Table~\ref{tab:covertype-complete-ranks} reports the complete rank comparison. Evidence uses preservation, contradiction uses actual inversion, and fragility uses endpoint-null alternate-context prediction change. The first two columns use 90 direction models; the last uses 45 fragility models. Parenthetical sample counts mark methods with partial coverage.

\begin{table*}[t]
\centering
\caption{Spearman correlation with realized behavior. SHAP interaction is available for random forests and XGBoost; KernelSHAP and LIME use their fixed formal subsets. $F$ values use the endpoint-null behavior target.}
\label{tab:covertype-complete-ranks}
\footnotesize
\setlength{\tabcolsep}{5pt}
\begin{tabular}{lrrr}
\toprule
Method & Preservation $\rho$ & Actual inversion $\rho$ & Endpoint-null change $\rho$ \\
\midrule
DECAF component & \textbf{0.864} & \textbf{0.987} & \textbf{0.974} \\
Endpoint $M$ & 0.657 & $-0.001$ & 0.148 \\
Abs & 0.804 & $-0.207$ & 0.588 \\
Signed net & \textbf{0.878} & $-0.353$ & 0.043 \\
SignFlip & $-0.919$ (78) & 0.888 (78) & 0.842 (39) \\
OppMass & $-0.412$ & 0.974 & 0.965 \\
Native SHAP & 0.781 & $-0.205$ & 0.481 \\
SHAP interaction & $-0.251$ (54) & 0.093 (54) & 0.069 (27) \\
KernelSHAP & 0.811 (30) & $-0.219$ (30) & 0.441 (15) \\
Global PFI & 0.664 & $-0.029$ & 0.676 \\
Context-conditioned PFI & 0.298 & 0.311 & 0.739 \\
PDP/ALE interaction & $-0.100$ & 0.827 & 0.937 \\
LIME & 0.845 (30) & $-0.384$ (30) & 0.506 (15) \\
\bottomrule
\end{tabular}
\end{table*}

For the two headline semantic coordinates, joint family/seed cluster bootstrap gives
\[
\rho(C,V_{\mathrm{inv}})=0.9868,
\qquad 95\%\ \mathrm{CI}=[0.9725,0.9967],
\]
and
\[
\rho(F,V_{\mathrm{null}})=0.9741,
\qquad 95\%\ \mathrm{CI}=[0.9416,0.9883].
\]
Evidence reaches $\rho(E,V_{\mathrm{keep}})=0.8642$ with 95\% interval $[0.8318,0.8955]$.

\subsection{The Same Treatment Produces Different Realized Mechanisms}
\label{app:covertype-family-audit}

Table~\ref{tab:covertype-family-audit} shows why treatment labels cannot serve as mechanism ground truth. The Invert generator produces genuine reversal in random forests and multilayer perceptrons, but mostly collapse in the remaining families. The Fragile generator becomes endpoint evidence in logistic regression, while nonlinear families retain substantial endpoint-null mass.

\begin{table*}[!tb]
\centering
\caption{Model-family audit. Direction columns average the Invert regime across two strengths and three seeds. Fragility columns average the Fragile regime across three seeds. ``Active'' is the endpoint-active fraction; ``Null change'' is the alternate-context prediction-change rate on endpoint-null examples.}
\label{tab:covertype-family-audit}
% \footnotesize
\resizebox{\textwidth}{!}{%
\begin{tabular}{lrrr|rrrrr}
\toprule
& \multicolumn{3}{c|}{Invert treatment} & \multicolumn{5}{c}{Fragile treatment} \\
Model family & $C$ & Invert rate & Collapse rate & $M$ & $E$ & $F$ & Active & Null change \\
\midrule
HistGradientBoosting & 0.000 & 0.000 & 1.000 & 0.020 & 0.070 & 0.180 & 0.340 & 0.623 \\
Logistic regression & 0.000 & 0.000 & 1.000 & 0.299 & 0.299 & 0.0001 & 0.989 & 0.000 \\
MLP & 0.088 & 0.743 & 0.257 & 0.028 & 0.084 & 0.080 & 0.350 & 0.219 \\
Random forest & 0.024 & 0.816 & 0.169 & 0.057 & 0.170 & 0.062 & 0.673 & 0.327 \\
XGBoost & 0.000 & 0.002 & 0.998 & 0.020 & 0.091 & 0.163 & 0.346 & 0.552 \\
\bottomrule
\end{tabular}%
}
\end{table*}

Within the Invert treatment alone, $C$ still tracks the amount of realized inversion with $\rho=0.922$. Among families with nonconstant inversion behavior, the correlations are 0.997 for histogram gradient boosting, 0.992 for random forests, 0.971 for XGBoost, and 0.940 for the MLP. Logistic regression has zero inversion throughout, so its within-family correlation is undefined.

\subsection{Fixed Semantic Readout and Magnitude Conditioning}
\label{app:covertype-fixed-semantic}

The fixed semantic benchmark assigns $E$, $C$, and $F$ to Evidence, Contradiction, and Fragility without fitting a meta-classifier. Table~\ref{tab:covertype-fixed-semantic} reports both the unconditional result and the within-Abs-bin result. The latter controls response magnitude through 12 quantile bins.

\begin{table*}[!tb]
\centering
\caption{Mechanism identification with fixed semantics and after conditioning on ordinary magnitude. Learned references use mechanism labels and are not access-matched to fixed DECAF.}
\label{tab:covertype-fixed-semantic}
\footnotesize
\begin{tabular}{lrrrr}
\toprule
Method & \makecell[r]{Fixed\\macro-AUROC} & \makecell[r]{Fixed\\macro-AUPRC} & \makecell[r]{Within-Abs\\macro-AUROC} & \makecell[r]{Within-Abs\\macro-AUPRC} \\
\midrule
\makecell[l]{DECAF / M / Abs /\\ Net / OppMass / SignFlip} & \textbf{0.816} & \textbf{0.768} & \textbf{0.956} & \textbf{0.948} \\
Endpoint $M$ & 0.487 & 0.437 & 0.565 & 0.666 \\
Abs & 0.518 & 0.450 & 0.495 & 0.585 \\
Native SHAP & 0.523 & 0.437 & 0.461 & 0.570 \\
SHAP interaction & 0.540 & 0.440 & -- & -- \\
\makecell[l]{Context-conditioned\\PFI} & 0.556 & 0.545 & 0.468 & 0.634 \\
PDP/ALE interaction & 0.539 & 0.438 & 0.559 & 0.739 \\
\addlinespace
\makecell[l]{Strong tabular reference\\(supervised)} & 0.987 & 0.968 & 0.983 & 0.994 \\
\makecell[l]{DECAF probe\\(supervised)} & 0.902 & 0.859 & 0.742 & 0.779 \\
\makecell[l]{Combined empirical\\reference (supervised)} & 0.985 & 0.965 & 0.983 & 0.994 \\
\bottomrule
\end{tabular}
\end{table*}

The fixed DECAF mechanism accuracy is 0.587. The supervised references use leave-one-model-family-out calibration, mechanism labels, and up to 15 input summaries. They bound cross-family decodability but do not replace the label-free comparison.

\subsection{Measured Cost and SHAP-Interaction Audit}
\label{app:covertype-cost}

Table~\ref{tab:covertype-cost-full} reports the registered cumulative worker time. Relative cost divides worker-seconds per model by the DECAF value. Coverage differs for methods whose formal budget uses a model subset.

\begin{table*}[htbp]
\centering
\caption{Complete measured method cost. SHAP interaction uses 128 stratified examples per tree model, split into four 32-example shards.}
\label{tab:covertype-cost-full}
\footnotesize
\begin{tabular}{lrrrr}
\toprule
Method & Models & Worker-s/model & Relative to DECAF & Predicted rows \\
\midrule
DECAF / $M$ / Abs / Net /  & 135 & 1.42 & $1.0\times$ & 25,920,000 \\
 OppMass / SignFlip &  &  &  &  \\
PDP/ALE interaction & 135 & 2.22 & $1.57\times$ & 38,880,000 \\
LIME & 45 & 8.38 & $5.92\times$ & 11,796,480 \\
PFI / context PFI & 135 & 13.32 & $9.41\times$ & 207,360,000 \\
KernelSHAP & 45 & 45.30 & $31.99\times$ & 377,501,760 \\
Native SHAP & 135 & 387.65 & $273.8\times$ & -- \\
Retraining reference & 135 & 19.70 & $13.91\times$ & -- \\
SHAP interaction & 54 & 15097.63 & $10662\times$ & -- \\
\bottomrule
\end{tabular}
\end{table*}

The formal SHAP-interaction run completed 216 shards over 54 tree models. It consumed at least 811,875 CPU-seconds (225.5 CPU-hours), used up to 32 concurrent shards, and required 17.71 hours of elapsed stage time. Its inversion correlation was only $0.093$ with a 95\% interval spanning zero.

\subsection{Matched-Pair Audit and Its Limits}
\label{app:covertype-matched-limit}

The automatic pair search produced only 11 protocol pairs. They cover 29.3\% of primary units, contain no same-family pair, and have a median relative endpoint-magnitude difference of 0.870 (maximum 0.945), although their Abs difference is small. DECAF reaches 0.727 mechanism accuracy on these pairs, but this set does not support a headline claim that both Abs and $M$ are matched. We therefore use the much better populated within-Abs-bin analysis in Table~\ref{tab:covertype-fixed-semantic} and retain the 11-pair result only as an audit.
